% !TeX root = ../thuthesis-example.tex

\chapter{引言}

\section{研究背景与意义}

\subsection{中性粒细胞的生物学功能及其在疾病中的作用}

中性粒细胞（neutrophil）是外周血中数量最为丰富的白细胞亚群，占循环白细胞总数的 50\%--70\%\cite{maierbegandt2024neutrophils}。在传统免疫学框架中，中性粒细胞被视为先天免疫防御的前线效应细胞，其主要通过趋化迁移、吞噬作用、颗粒酶释放以及还原型烟酰胺腺嘌呤二核苷酸磷酸（reduced nicotinamide adenine dinucleotide phosphate, NADPH）氧化酶复合体介导的活性氧（reactive oxygen species, ROS）爆发等机制清除入侵病原体，从而快速限制感染扩散\cite{mayadas2014multifaceted}。然而，近年来随着单细胞转录组学、空间组学及体内成像技术的发展，中性粒细胞的功能认知已发生显著转变。越来越多的研究表明，其并非单一均质群体，而是在分化阶段、组织微环境及病理背景影响下呈现高度动态的功能状态多样性\cite{arocacrevilln2024neutrophils}。

在生理条件下，中性粒细胞除参与抗感染反应外，还在组织修复、血管稳态维持及细胞外基质（extracellular matrix, ECM）重塑过程中发挥调节作用。例如，其释放的基质金属蛋白酶和弹性蛋白酶可调控局部炎症微环境并影响血管生成过程\cite{rosales2018neutrophil}。在病理条件下，中性粒细胞与多种重大疾病的发生发展密切相关。特别是在心血管疾病中，中性粒细胞胞外诱捕网（neutrophil extracellular traps, NETs）被证实参与动脉粥样硬化进展及血栓形成，是无菌性炎症放大的关键分子事件之一\cite{baaten2024netosis}。此外，在肿瘤微环境和自身免疫疾病中，中性粒细胞亦表现出功能重编程现象，对疾病进程产生重要影响。

综上，中性粒细胞已从传统意义上的急性杀伤细胞，转变为连接感染防御、组织修复与慢性炎症调控的重要免疫枢纽。其功能执行依赖复杂分子网络的精细调控，这为探究其特异性分子依赖模式提供了理论基础。

\subsection{蛋白质必需性理论与功能基因组学研究进展}

蛋白质必需性（protein essentiality）是生命科学与系统生物学研究中的核心概念，通常指编码该蛋白的基因在特定生物体系中对生物体的生存、发育及关键生理功能维持具有不可替代性\cite{zhang2016predicting,xiao2015identifying}。必需蛋白作为细胞功能网络中的关键分子节点，广泛参与基础代谢过程、能量稳态维持、信号转导调控及细胞周期控制等核心生命活动，其功能缺失往往导致细胞死亡、发育异常或个体功能障碍\cite{jeong2001lethality}。

在功能基因组学层面，蛋白质必需性通常通过基因扰动后细胞适应度（fitness）变化进行定量评估。随着成簇规律间隔短回文重复序列及其相关蛋白 9 技术（clustered regularly interspaced short palindromic repeats-associated protein 9, CRISPR-Cas9）的发展，研究者可构建包含大量单导 RNA（single-guide RNA, sgRNA）的混合文库，在细胞群体中进行竞争性生长筛选，通过测序分析 sgRNA 丰度变化来推断基因对细胞存活的贡献\cite{shalem2014genome}。为提高筛选结果的可靠性，研究者发展了多种统计分析与误差校正方法，例如用于 sgRNA 富集分析的基于模型的全基因组 CRISPR-Cas9 敲除分析方法（Model-based Analysis of Genome-wide CRISPR-Cas9 Knockout, MAGeCK）算法，以及用于校正拷贝数变异影响的 CERES 拷贝数效应校正模型\cite{meyers2017computational}。这些方法的建立使得在全基因组尺度上系统识别必需基因进而筛选必需蛋白成为可能。

根据功能重要性的广泛程度，必需蛋白可进一步划分为“核心必需（core essential）”与“情境特异必需（context-specific essential）”。核心必需蛋白在多种细胞类型中均对细胞存活与基本生命活动具有关键作用，反映其在维持普遍生理过程中的基础地位；而情境特异必需蛋白则仅在特定遗传背景、分化阶段或功能状态下发挥不可替代作用。这强调，蛋白质的重要性并非固定不变，而是受到细胞类型与生理状态的调控\cite{tsherniak2017defining}。

从分子网络角度来看，情境特异必需蛋白往往参与某一特定细胞状态下被强化或激活的功能通路。例如，在高代谢负荷、强氧化应激或快速增殖条件下，某些蛋白可能成为维持细胞稳态的关键节点，而在其他状态下则并非决定性因素。因此，必需蛋白不仅体现生命活动的基本构成，还反映不同细胞在特定生理或病理环境中的功能侧重。

这为精准医学提供了重要理论基础。若某一细胞类型或病理状态必须依赖特定蛋白维持其异常或过度活跃的功能过程，则该蛋白可能成为具有选择性的干预靶点。相较于在多种组织中均发挥基础功能的核心必需蛋白，细胞类型或状态特异的必需蛋白更可能在实现治疗效果的同时降低系统性毒副作用\cite{hart2015high}。因此，围绕情境特异必需蛋白构建分子干预策略，已成为药物研发与治疗优化的重要方向。

综上，蛋白质必需性的研究不仅有助于揭示生命系统的基本运行规律，也为筛选具有选择性优势的治疗靶点提供理论依据。然而，在免疫细胞体系中，特别是在执行高强度炎症与应激功能的中性粒细胞中，其特异性必需蛋白谱及关键调控分子仍有待系统阐明。

\subsection{细胞类型与生理背景对蛋白质必需性的影响}

尽管蛋白质必需性常用于描述某一基因对细胞存活的重要程度，但越来越多的研究表明，必需蛋白并非一个固定不变的分子集合，而是受到细胞类型、分化阶段及生理背景的显著影响。不同细胞由于发育来源、功能定位及代谢特征存在差异，其维持基本稳态所需的关键分子亦不完全相同。这一现象在多种人源细胞系的系统性功能研究中得到验证，不同细胞类型之间的必需基因组成存在明显差异\cite{blomen2015gene,wang2015identification}。

进一步的分析显示，基因的重要性不仅取决于细胞类型本身，还受到功能状态与外界环境的调控。例如，在营养限制、氧化应激增强或信号通路持续激活等条件下，部分分子在维持细胞稳态中的作用被显著强化，而在静息或低负荷条件下则不再表现为关键节点\cite{birsoy2014metabolic}。这种由生理背景驱动的功能重要性变化，说明蛋白质必需性具有明确的背景依赖特征。

此外，细胞分化与功能重编程亦会改变分子需求结构。单细胞层面的功能研究指出，不同分化阶段或功能亚群之间在关键调控分子的使用上存在显著差异\cite{replogle2022mapping,zhou2023single}。在免疫系统中，细胞活化与效应功能执行往往伴随代谢重塑与信号网络重构，从而改变特定分子的功能地位\cite{xie2020single}。因此，必需蛋白的界定需要结合具体细胞功能状态进行分析，而不能简单应用到其他细胞体系。

对于中性粒细胞而言，这种细胞类型与功能状态相关的必需性差异尤为突出。中性粒细胞在执行吞噬作用、ROS 爆发及NETs 形成等功能时，处于高度代谢与氧化压力环境，其生存与功能维持依赖特定分子通路的精细调控。然而，目前公开的大规模必需基因数据主要来源于肿瘤细胞系，对终末分化且功能高度特化的免疫细胞覆盖相对有限。因此，在中性粒细胞这一特定细胞类型中系统识别并解析其特异性必需蛋白，对于理解其在炎症与免疫反应中的调控机制具有重要科学意义。

综上所述，蛋白质必需性呈现出显著的细胞类型特异性和生理背景依赖性。这一特征为针对特定细胞类型开展必需蛋白系统研究提供了理论基础，也为后续构建中性粒细胞特异必需蛋白预测框架奠定了生物学前提。

\subsection{基于 PLM 的必需性预测方法及研究意义}

在特定细胞类型中系统识别必需蛋白，往往面临现实实验条件的限制。以中性粒细胞为例，该类细胞属于终末分化细胞，寿命较短，体外培养时间有限，难以开展长期适应度评估实验。同时，成熟中性粒细胞体外扩增能力较弱，基因编辑效率较低，使得大规模功能筛选在技术上具有较高难度。此外，中性粒细胞在执行吞噬、ROS 爆发及NETs 形成等功能时呈现高度动态变化，也增加了实验定量分析的复杂性。因此，仅依赖实验手段难以全面解析中性粒细胞特异性必需蛋白谱。

在此背景下，计算方法成为拓展必需蛋白研究的重要补充途径。早期研究表明，蛋白质相互作用（protein--protein interaction, PPI）网络中的拓扑位置与基因重要性存在关联\cite{jeong2001lethality}，基于网络结构与表达特征的建模方法在一定程度上提高了预测效率。然而，这类方法通常依赖人工特征构建，对数据完整性与网络质量要求较高，在跨细胞类型场景中的泛化能力有限。

近年来，随着深度学习与自监督学习技术的发展，蛋白语言模型（protein language model, PLM）为从氨基酸序列中直接提取功能表征提供了新的技术路径。大规模自监督训练表明，模型能够从海量序列中学习到反映结构与进化约束的隐含表示\cite{rives2021biological}。进一步的研究显示，扩展模型规模与训练数据后，其内部表示可用于结构预测与功能推断等多种下游任务\cite{lin2023evolutionary}。这种序列驱动的表征方式在实验数据受限的情况下具有重要潜力。

因此，在实验筛选存在客观限制的前提下，引入基于 PLM 的计算预测方法，为在特定细胞类型背景下识别潜在关键靶点蛋白提供了可行路径。本研究正是在这一方法学背景下，构建面向中性粒细胞的必需蛋白预测框架。

\section{国内外研究现状}

\subsection{传统蛋白质必需性鉴定方法及其局限性}

蛋白质必需性的界定起源于经典遗传学研究，其核心逻辑是通过基因功能缺失所导致的生存或发育障碍来推断关键分子集合。早期在酿酒酵母中的系统性单基因缺失研究首次构建了全基因组尺度的必需基因图谱\cite{giaever2002functional}，并进一步揭示必需基因在 PPI 网络中往往处于拓扑中心位置。此类研究奠定了“中心性—必需性”之间的统计关联框架，也推动了功能基因组学向系统层面发展。此外，最小基因组研究从系统生物学角度进一步强化了“核心生命功能集合”的概念。Hutchison 等通过转座子插入突变系统分析支原体基因组，提出最小生存基因集合（minimal gene set）的框架，强调维持细胞基本生命活动所需的核心基因模块\cite{hutchison1999global}。随后，通过人工合成最小细胞基因组的研究进一步验证，在严格控制的培养条件下，可以通过去除非必需基因构建可稳定存活的简化基因组体系\cite{hutchison2016design}。这些研究从理论层面巩固了“核心必需基因”作为维持生命基本结构与功能的稳定模块这一认识，并推动了将必需性视为相对稳定生物学属性的研究范式。然而，这些研究通常基于稳定遗传背景与标准培养条件，其所得结论隐含“基因功能重要性具有跨条件稳定性”的假设。

随着研究体系从模式生物拓展至哺乳动物细胞，RNA 干扰（RNA interference, RNAi）技术为大规模功能筛选提供了可行路径\cite{mohr2014rnai}。RNAi 通过抑制 mRNA 表达实现功能削弱，在多种细胞模型中用于识别存活相关基因。然而，大量研究指出 RNAi 存在脱靶效应与非特异性毒性问题\cite{jackson2010recognizing}，且不同 siRNA 的敲降效率差异较大，使结果在不同实验平台间难以直接比较。进一步的系统评估显示，不同 RNAi 文库之间的筛选结果重叠度显著低于理论预期，提示技术噪声与脱靶效应对必需基因判定具有实质影响\cite{echeverri2006minimizing,birmingham2006utr}。然而，相较于技术层面的偏差，更为关键的是 RNAi 筛选所依赖的“细胞适应度下降”判定框架。在多数研究中，基因必需性被等同于敲降后细胞增殖能力或存活率的显著下降\cite{hart2015high}。这种定义方式在高度增殖的肿瘤细胞系中具有可操作性，但其本质反映的是特定培养条件下的生长依赖性，而非广义生理环境中的功能必需性。因此，RNAi 阶段所构建的必需基因集合在跨细胞类型、跨功能状态的推广过程中，天然受到“增殖依赖范式”的结构性限制。此外，RNAi 并非完全失活策略，其对低表达或功能冗余基因的干扰效果有限，从而可能低估某些分子的真实功能重要性。这些技术特征使得 RNAi 所构建的必需基因集合在精度与可重复性方面受到质疑。

CRISPR-Cas9 技术的出现显著提升了基因组编辑的精确性与通量，使得在细胞层面开展全基因组尺度的功能筛选成为现实。基于混合 sgRNA 文库的竞争性筛选策略能够通过高通量测序定量分析扰动后 sgRNA 的丰度变化，从而推断基因对细胞适应度的贡献。后续研究通过改进 sgRNA 设计原则\cite{doench2016optimized}与统计建模方法，进一步提高了筛选数据的可靠性与特异性。大规模项目（如癌症依赖图谱项目，Cancer Dependency Map, DepMap）构建了跨数百细胞系的基因依赖图谱，为系统解析肿瘤细胞的功能脆弱性提供了资源基础。除单基因必需性之外，近年来研究进一步揭示了基因功能在网络层面的复杂性。值得注意的是，随着大规模 CRISPR 数据的积累，对筛选结果的再分析逐渐揭示必需基因集合在不同细胞类型之间的显著差异。系统比较 300 余种癌细胞系的依赖图谱发现，仅有少数基因在多数细胞系中表现为“核心必需”，而大量基因呈现明显的细胞类型特异性依赖\cite{behan2019prioritization}。进一步的统计模型分析提出，应区分“核心必需基因”与“情境特异必需基因”，后者的必需性往往取决于特定代谢状态、信号通路激活情况或基因组背景\cite{hart2016bagel}。

尽管如此，CRISPR 筛选体系仍依赖若干关键前提：其一，细胞能够稳定培养并维持足够时间以观察适应度差异；其二，基因组结构相对稳定，以避免拷贝数变异对切割信号的干扰；其三，细胞群体间竞争模型能够准确反映单基因扰动对生长优势的影响。然而，在非增殖型或终末分化细胞中，上述前提往往难以满足。以免疫细胞为例，原代 T 细胞或巨噬细胞的基因编辑效率与体外维持时间均显著低于肿瘤细胞系\cite{shifrut2018genome}。与可持续增殖的肿瘤细胞系不同，原代中性粒细胞属于终末分化细胞，生命周期短、转染效率低且对基因编辑高度敏感，使得大规模 CRISPR 筛选在该细胞类型中难以实施。即便在诱导分化模型中，编辑效率与筛选窗口期仍存在显著限制。因此，在中性粒细胞这一高度动态、应激响应依赖型的免疫细胞中，传统功能基因筛选技术难以构建完整的必需性图谱。现有研究多通过祖细胞分化模型或 HL-60 人早幼粒细胞白血病细胞系（human promyelocytic leukemia cell line, HL-60）等替代细胞体系进行间接推断\cite{nasri2020crispr}。这种基于细胞系分化模型的间接推断，实质上是一种牺牲“生物学保真度”以换取“遗传可操作性”的折衷方案。HL-60 等替代体系在诱导分化后，其表观遗传景观、蛋白质水解酶组分以及颗粒蛋白的成熟度，均与外周血原代中性粒细胞存在显著的生化差异 \cite{allen2023closing}。这意味着，在替代体系中鉴定出的“必需蛋白”，往往无法真实还原原代细胞在炎症反应中的功能依赖路径。如何在不破坏细胞稳态的前提下，直接在具备完全生物学活性的原代终末分化细胞中，构建高通量的蛋白质功能扰动与因果验证平台，是当前功能基因组学亟待跨越的边界。

综上所述，传统实验方法在因果验证与机制解析方面具有不可替代的优势，但其可扩展性与跨细胞类型泛化能力受到实验体系与生物学条件的双重限制。尤其在功能高度动态、状态快速变化的免疫细胞系统中，如何构建既具备生物学真实性又具有跨背景可迁移性的必需蛋白识别框架，仍是当前领域亟待解决的问题。所以基于蛋白序列信息构建的深度学习预测框架不仅是实验方法的补充工具，更成为突破实验可行性边界的关键技术路径。通过对大规模序列语义空间的系统建模，可在无需直接基因扰动的前提下，从全基因组层面推断潜在功能依赖模式，为原代免疫细胞中的关键调控蛋白筛选提供理论依据。

\subsection{早期计算预测方法：网络与浅层序列特征驱动}

在深度表示学习方法引入之前，必需蛋白的计算预测主要依赖两类范式：一类以PPI 网络（protein--protein interaction, PPI）拓扑特征为核心，另一类以蛋白序列及功能注释特征为输入构建传统机器学习分类器。二者分别从系统结构与分子属性层面建模蛋白重要性，其方法学演进过程构成了必需蛋白预测早期阶段的主要技术框架。

网络驱动预测方法的理论基础来源于 Jeong 等发表于Nature的研究\cite{jeong2001lethality}。该研究基于酿酒酵母的 PPI 网络数据，将节点度（degree centrality）与基因敲除实验获得的必需基因集合进行统计关联分析，发现高度连接节点中必需基因比例显著升高。作者据此提出“中心性–致死性规则”（centrality–lethality rule），为后续将网络拓扑指标纳入预测模型提供理论依据。随后Hahn 与 Kern\cite{hahn2005comparative}在酵母线虫与果蝇三个真核物种中构建跨物种比较框架，对度中心性、介数中心性与紧密中心性等拓扑指标进行了系统统计分析。结果表明，多种中心性指标在不同物种中均与基因必需性显著相关，且关联方向一致。这一研究在跨物种层面强化了“网络位置约束功能重要性”的理论框架，使中心性从单一统计相关上升为具有潜在机制含义的结构特征。

然而，He 与 Zhang\cite{he2006why}的研究对“中心性–致死性规则”提出了反驳。他们指出，枢纽蛋白更可能为必需蛋白，并不必然源于其在维持网络整体结构中的关键作用，而可能仅仅由于其参与的相互作用数量更多，从而更高概率地涉及“必需相互作用”（essential PPIs）。在其构建的概率模型中，蛋白必需性的发生可由简单的组合概率解释，而无需引入网络拓扑结构的全局性质。进一步分析表明，在控制度中心性后，介数中心性与紧密中心性等全局中心性指标不再表现出额外解释力。这表明 PPI 网络本质上是静态、汇总式的结构图谱，其边的存在并不区分相互作用的条件依赖性、特异性或调控状态。在此框架下，将必需性视为网络拓扑位置的函数，实质上忽略了蛋白功能的动态调控本质。换言之，中心性指标所捕捉的更多是“连接频度”的统计特征，而非蛋白在特定生理或应激情境下的功能依赖机制。

在拓扑特征改进方面，研究人员对网络模块结构与必需基因分布进行了再分析，指出必需基因更倾向于位于高保守功能模块内部\cite{yu2008genomic}。其方法通过识别网络模块并计算模块内部基因必需性富集程度，说明局部网络结构信息可能比全局度指标更具解释力。该工作推动预测特征从“单节点中心性”向“模块化组织特征”扩展。然而，其分析仍基于静态整合网络，未考虑不同条件下模块活跃度变化。

关于网络静态性的局限的研究同样也有进展，研究人员整合表达数据与互作网络，发现不同生理条件下活跃的功能模块发生切换，提出互作网络具有动态模块化特征\cite{han2004evidence}。尽管该研究并非直接针对必需性预测，但其结果从系统层面表明：静态整合网络仅是多状态网络的平均表示。因此，基于静态拓扑构建的预测模型难以刻画细胞情境依赖的功能脆弱性。

在人类互作网络层面，Luck 等[5]构建的高置信度二元互作图谱（HuRI）系统评估了现有互作数据库的覆盖率与偏倚\cite{luck2020reference}。该研究通过统一实验平台获得约50000条二元互作数据，并指出不同实验体系之间存在显著差异。其方法学意义在于揭示当前 PPI 网络并非完整生物系统的精确表示，而是实验条件下的观测子集。这意味着，基于现有 PPI 网络构建的预测模型，其特征输入本身存在不确定性来源。

与网络拓扑方法并行发展的，是基于蛋白序列与功能注释的预测路径。该类方法的共同前提是：蛋白必需性能够通过若干可解释的统计特征或功能标签加以刻画。这一前提在早期数据规模有限、功能认知尚不完善的背景下具有合理性，但随着数据规模与问题复杂性的提升，其内在张力逐渐显现。

早期研究主要围绕氨基酸组成、二肽频率及理化性质参数展开。这类特征具有明确生物学含义，能够解释为“组成偏好”或“结构稳定性倾向”。然而，其表达能力受限于一阶统计结构。即便引入进化保守性等序列特征以增强表示维度 \cite{deng2011investigating}，本质上仍然是对局部替代概率的线性展开，而非对序列全局依赖结构的建模。以位置特异性评分矩阵（position-specific scoring matrix, PSSM）为代表的保守性表示在部分场景中确实有效，提示进化约束与必需性存在关联；但其局限同样明显：当同源信息不足或蛋白进化速率较快时，该类表示会快速退化。

在算法层面，研究者逐渐意识到类别不平衡与特征冗余对性能的影响，因而引入平衡随机森林与特征选择策略 \cite{anaissi2013balanced}。这些改进提升了模型稳定性，却并未改变表示本身的能力边界。性能提升来源于决策函数的优化，而非对蛋白表示方式的根本重构。

随后，信息论特征被纳入框架 \cite{nigatu2017sequence}。序列复杂度与信息熵被视为潜在区分信号。这一尝试体现出对“统计频率不足”的反思，但信息熵仍属于低阶统计量，无法表达残基间非线性耦合，也无法反映三维结构约束。因此，其改进效果呈现明显的边际递减趋势。

功能注释的引入在一定程度上改变了问题结构。通过整合基因本体论（Gene Ontology, GO）语义信息 \cite{zhang2018detecting}，模型开始显式利用已有生物学知识。这种策略增强了可解释性，却带来了新的依赖：预测能力受限于注释完整性。注释本身是研究积累的结果，而非蛋白本体属性。对于未充分研究的蛋白，或特定细胞类型中与情境相关的依赖特征，该类方法难以提供有效预测。

从整体上看，这一阶段的方法演进呈现出“特征扩展—性能提升—边际收敛”的轨迹。其根本原因并非序列信息本身不足，而在于对序列的表示方式长期停留在人工构建的低阶统计特征层面。无论是氨基酸组成、多序列比对衍生的保守性矩阵（如 PSSM），还是信息熵与理化参数，本质上均属于预定义的特征压缩方式。这种方式假设蛋白功能可以通过有限维、可解释的统计量充分表达。然而，蛋白质序列本身蕴含着复杂的高阶依赖结构，包括远距离残基协同、进化耦合关系以及潜在的结构约束模式。这些信息虽然存在于原始序列中，但难以通过人工统计特征显式捕获。因此，传统方法的性能瓶颈更应理解为表示学习能力的瓶颈，而非序列信息量的瓶颈。

因此，问题的关键并不在于是否引入更多外部信息，而在于：是否能够在不增加额外模态数据的前提下，从原始序列中自动学习更具表达力的高阶表示。当研究焦点从“特征构建”转向“表示学习”时，基于自监督预训练的 PLM 应运而生 \cite{rives2021biological,lin2023evolutionary}。与传统人工特征不同，PLM 直接在大规模序列数据上学习隐含的进化与结构模式，使序列表示从低阶统计描述升级为高维语义嵌入表示。这一转变并未改变信息来源（仍为序列），而是改变了信息利用方式，从而在表示层面突破了传统方法的结构性瓶颈。

综上所述，在深度表示学习方法引入之前，必需蛋白预测主要沿两条技术路径展开：其一基于 PPI 网络拓扑结构，通过中心性或模块化特征推断功能重要性；其二基于序列组成与功能注释，依赖人工构建的统计特征结合传统机器学习分类器进行判别。两类方法分别从系统结构与分子属性层面对蛋白必需性进行了建模，在数据规模有限、特征工程主导的阶段具有重要推动作用。网络驱动模型依赖于静态整合的 PPI 图谱，难以刻画条件依赖性与细胞情境差异；序列驱动模型虽摆脱了网络数据依赖，却长期停留在低阶统计特征表达框架，无法充分挖掘蛋白序列中隐含的远距离残基耦合与高阶进化约束模式。正是在这一背景下，基于自监督预训练的 PLM 逐渐成为新的技术转折点，为必需蛋白预测从“特征工程驱动”向“表示学习驱动”转型提供了理论与方法基础。

\subsection{深度学习与 PLM 的兴起}

蛋白质序列到结构与功能的映射关系具有显著的多尺度复杂性。传统计算方法通常基于人工构建的特征体系，如氨基酸组成统计、理化性质指标、多序列比对衍生的进化保守性评分及二级结构预测结果等，并结合支持向量机或随机森林等浅层模型进行建模。然而，此类方法的表达能力受限于特征设计的先验假设，难以系统刻画序列中的长程依赖关系与高阶非线性结构。深度学习方法的引入，从方法论层面突破了这一表达瓶颈。

在深度学习介入初期，研究重点在于利用深度神经网络自动学习序列表征，以替代传统依赖于专家经验的人工特征工程（如氨基酸组成、自相关系数等）。
卷积神经网络（convolutional neural network, CNN）通过多层局部卷积滤波器提取序列片段中的模式特征，在蛋白质二级结构预测 \cite{wang2016protein} 以及翻译后修饰位点识别 \cite{wang2017musitedeep} 等任务中展现出较强的局部特征建模能力。
由于卷积算子的局部感受野限制，传统 CNN 在缺乏多层堆叠或扩张卷积机制的情况下，难以充分整合蛋白序列中的长程依赖关系。为弥补这一不足，循环神经网络（recurrent neural network, RNN）及其变体长短期记忆网络（long short-term memory, LSTM）被引入序列建模领域 \cite{hochreiter1997long}，通过递归结构在时间维度上传递隐藏状态，从而在一定程度上缓解长程依赖建模中的梯度消失问题。随后，多项研究将双向 LSTM 应用于蛋白结构与功能预测任务，验证其在长序列建模中的优势。
尽管监督学习方法在特定物种数据集上能够取得较高预测精度，但其性能高度依赖充足且分布一致的标注数据。必需基因的实验鉴定通常依赖饱和转座子诱变（Tn-Seq）等高通量技术 \cite{vanopijnen2013transposon}，该类实验流程复杂、成本高昂，且现有数据主要集中于少数模式生物 \cite{zhang2004deg,hutchison2016design}。在标注数据匮乏或物种分布发生变化的情况下，监督模型易受到过拟合与分布偏移问题的影响，从而导致泛化性能显著下降。

随着结构生物学的发展，研究者意识到蛋白质的功能及其必需性不仅取决于线性序列，更取决于其三维空间拓扑以及在复杂PPI中的交互地位。
图神经网络（Graph Neural Networks, GNN）的引入为蛋白质建模提供了新的归纳偏置。通过将残基视为节点，将空间接触或功能关联视为边，GNN 利用消息传递机制实现节点表征的动态更新。其核心算子通常表达为：
\begin{equation}
h_i^{(l+1)} = \sigma \left( \sum_{j \in \mathcal{N}(i)} \alpha_{ij} W^{(l)} h_j^{(l)} \right)
\end{equation}
其中，$\alpha_{ij}$ 代表注意力权重或聚合系数\cite{gilmer2017neural,kipf2017semi}。
这种建模方式成功捕捉了必需蛋白在代谢网络中的“中心性-致死性”特征，即处于网络关键枢纽位置的蛋白更倾向于具有必需性。

然而，显式结构建模仍面临现实约束。一方面，实验解析的蛋白质三维结构数量远低于可获得的序列规模 \cite{consortium2019protein}，导致基于结构监督的模型训练数据受限；另一方面，传统基于静态 PPI 网络的图模型通常忽略相互作用的时空条件依赖性，而大量研究表明蛋白质互作网络具有显著的动态模块化特征 \cite{han2004evidence}。这种静态近似限制了模型在复杂生理环境下识别情境特异必需蛋白的能力。

随着自然语言处理领域表示学习思想的发展，研究者开始探索蛋白序列是否能够通过自监督方式学习统一的嵌入表示。早期工作如 Heinzinger 等\cite{heinzinger2019modeling}采用类 ELMo 架构在大规模蛋白序列上进行语言模型预训练，并验证其特征嵌入在二级结构与定位预测任务中的迁移能力，首次提出“蛋白序列语言化建模”的可行性。随后，Bepler 与 Berger\cite{bepler2019learning} 将结构接触信息融入序列嵌入特征学习框架，通过联合监督与无监督信号优化表示空间，使可训练的嵌入特征同时编码序列与结构约束，为后续结构感知型表示学习奠定基础。

在此基础上，Alley\cite{alley2019unified} 等提出 UniRep 模型，在超过 2000 万条未标注蛋白序列上训练循环语言模型，并将隐藏状态作为统一序列嵌入。实验结果表明，该 embedding 在突变效应预测与蛋白工程优化任务中可直接迁移使用，且无需复杂特征工程。这一研究系统性证明蛋白序列可采用语言建模范式进行通用表示学习。

随后，Rao 等首先构建蛋白嵌入任务评估基准（Tasks Assessing Protein Embeddings, TAPE）框架，对多种预训练策略在二级结构预测、接触预测与稳定性评估等任务上的迁移性能进行系统比较，结果表明预训练显著提升下游表现，并揭示模型容量与训练规模对性能具有重要影响 \cite{rao2019evaluating}。然后进一步分析 Transformer PLM 内部注意力权重与真实残基接触图之间的相关性，发现部分注意力头能够显式对应空间邻近残基，表明预训练模型在无监督条件下已学习到结构相关信息 \cite{rao2021transformer}。

这一阶段的研究完成了从“任务驱动建模”向“统一表示学习”的方法学转变，使蛋白序列嵌入成为后续大规模预训练模型发展的理论基础。
Transformer 架构的引入为蛋白序列建模提供了关键技术基础\cite{vaswani2017attention}。与循环神经网络相比，自注意力机制能够在单层计算中直接建立任意位置残基之间的依赖关系，从而有效建模蛋白序列中普遍存在的长程相互作用。考虑到三维结构中空间邻近残基往往在序列上相距较远，Transformer 的全局注意力结构天然契合蛋白折叠问题的依赖模式。

在结构预测层面，AlphaFold2 通过引入 Evoformer 模块，将多序列比对（multiple sequence alignment, MSA）表示与基于注意力机制的残基对表示进行迭代交互更新，实现进化共变信息与空间约束的深度融合，并在第14届蛋白质结构预测关键评测（Critical Assessment of Structure Prediction, 14th edition, CASP14）中达到接近实验精度的三维结构预测水平 \cite{jumper2021highly}。尽管该模型仍依赖多序列比对作为输入，其成功从任务层面验证了注意力机制在整合跨残基依赖关系方面的有效性。结合前述无监督语言模型中结构信息的涌现现象，可以认为，基于注意力机制的深度架构为蛋白序列向结构映射提供了统一且可扩展的技术基础。

在此背景下，Rives 等构建进化尺度建模蛋白语言模型（Evolutionary Scale Modeling, ESM）中的 ESM-1b 模型，在超过 2.5 亿条 UniRef50 蛋白序列上进行掩码语言模型（masked language modeling, MLM）预训练\cite{rives2021biological}。该模型包含数亿级参数，通过随机遮蔽部分氨基酸残基并预测其类别，使模型被迫利用上下文信息重建局部序列，从而在无监督条件下学习序列内部的统计规律。不同于早期基于监督标签的模型，ESM-1b 的训练目标完全独立于特定下游任务，因而其隐藏表示具有更强的通用性。

更为重要的是，Rives 等对模型内部表示进行了系统分析。他们发现：（1）部分注意力头的权重分布与实验结构数据中的残基接触图高度相关；（2）隐藏层 embedding 可用于线性预测二级结构类别；（3）模型对特定保守位点的预测置信度与进化保守性评分呈显著相关关系。这些结果表明，模型在未显式接收结构或功能监督信号的情况下，自发学习到与空间结构和进化约束相关的统计特征。由此，蛋白序列的结构与功能信息被证明可以通过规模化自监督训练自然编码于潜在表示空间之中。

在 ESM 路线之外，ProtTrans 系列模型从架构与训练策略层面对 PLM 进行了系统扩展\cite{elnaggar2022prottrans}。Elnaggar 等采用双向编码器表征 Transformer（Bidirectional Encoder Representations from Transformers, BERT）、XLNet 与文本到文本迁移 Transformer（Text-to-Text Transfer Transformer, T5）等多种 Transformer 架构，在 UniRef50 及更大规模数据集上进行预训练，并构建参数规模从百万级到十亿级的多个模型版本。与 ESM 侧重结构信息验证不同，ProtTrans 更强调迁移学习能力的系统评估。研究者在亚细胞定位、二级结构预测、跨膜蛋白识别等多项任务上进行基准测试，结果显示预训练模型在无需复杂特征工程的情况下即可显著提升性能，且模型规模与训练数据规模的增加带来持续增益。这一结果从经验层面强化了“规模驱动规律”在蛋白建模中的适用性。

ESM 与 ProtTrans 的互补性在于：前者通过结构层分析证明预训练表示具有生物物理含量，后者通过多任务评估验证其迁移能力与泛化性能。两条路线共同构建了 PLM 的理论与实证基础。

由此，PLM 逐渐从“方法探索”演化为“统一建模范式”。其核心特征包括：
（1）在大规模未标注序列上进行自监督预训练；
（2）通过掩码预测任务学习上下文敏感表示；
（3）在潜在空间中整合结构约束、进化保守性与功能信号；
（4）通过下游微调或特征提取实现跨任务迁移。
相较于早期任务特定型深度模型，PLM 不再针对单一生物学问题进行结构设计，而是提供一个统一的高维表示空间，使不同层级的蛋白质问题可以在同一表征框架内进行建模与比较。这一方法学转变标志着蛋白序列建模正式进入“预训练基础模型”阶段。

Meier 等利用预训练 PLM 进行零样本突变效应预测，发现突变前后序列的对数似然差值与实验测得的功能改变显著相关，表明模型学习到的概率分布能够反映进化选择压力与功能约束强度 \cite{meier2021language}。ProteinBERT 通过联合序列建模与功能注释任务进行多任务预训练，在多项 GO 功能预测任务中表现出良好的迁移能力，进一步验证了预训练语言模型在功能注释领域的泛化潜力 \cite{brandes2022proteinbert}。

综上所述，蛋白建模方法经历了由人工特征驱动向端到端深度学习、再向大规模自监督预训练与统一表示空间构建的演进过程。PLM 的形成不仅提升了下游预测任务的性能上限，更从方法论层面确立了“预训练基础模型 + 下游任务适配”的统一技术框架，为结构、功能及进化层面问题的联合建模提供了共享表示空间。

蛋白质序列作为一种蕴含高度复杂生物信息的“生命语言”，其计算表征的范式演迁构成了过去十年生物信息学研究的主旋律。从最初对理化特征的端到端监督学习，到对空间拓扑结构的显式建模，直至今日以大规模自回归或自编码预训练为核心的PLM，这一演进过程不仅显著提升了预测精度，更深刻地改变了人类对基因必需性等复杂生物学属性的建模维度。

在必需蛋白预测任务中，PLM 的引入具有直接方法学意义：必需性本质上与蛋白的结构稳定性、功能模块化与进化约束密切相关，而这些信息在序列中往往以复杂的统计规律隐式存在。PLM 通过大规模预训练学习到的序列表示，能够在一定程度上缓解传统浅层序列特征表达能力不足的问题，为必需蛋白预测提供新的性能上限空间。

\subsection{PIC模型及其方法学框架}

在基于 PLM 的必需蛋白预测研究中，PIC（Protein Importance Calculator）是目前具有代表性的序列驱动方法之一。该模型由 Kang 等发表于 Nature Computational Science（2024），系统性验证了大规模预训练 PLM 在蛋白必需性识别任务中的应用潜力 PIC 模型\cite{kang2024comprehensive}。

在方法框架上，PIC 以预训练 PLM 作为序列表征模块，通过微调策略构建端到端分类模型，并在统一架构下实现跨层级预测体系。具体而言，PIC 构建了人类层面模型（PIC-human）、小鼠层面模型（PIC-mouse）以及 323 个细胞系特异模型（PIC-cell），并采用软投票（soft voting）集成策略整合细胞系预测结果，从而形成细胞层面的综合评估框架。该设计实现了从“单一物种预测”向“跨物种—跨细胞层级预测”的拓展，为蛋白必需性在不同生物学层面上的系统评估提供了方法学基础。

在序列表征层面，PIC 通过系统消融实验比较多种特征提取与池化策略，最终选用参数规模为 6.5 亿的 ESM-2 作为 backbone 模型，并结合平均池化生成蛋白级表示。在独立测试集上，PIC-human 的样本特征曲线下面积（area under the receiver operating characteristic curve, AUROC）达到 0.9132，PIC-mouse 达到 0.8736；在细胞系层面，与 DeepCellEss 等既有方法相比，PIC 在 323 个细胞系上的 AUROC 与精准率-召回率曲线下面积（area under the precision-recall curve, AUPRC）平均提升分别为 9.64\% 与 10.52\%，体现出预训练序列表示在必需性预测任务中的显著优势。

此外，PIC 基于模型输出概率构建蛋白必需性评分（protein essential score, PES），将传统二分类任务拓展为连续量化评估框架，并通过进化保守性、表达水平与功能富集分析验证其生物学一致性。这一设计为必需蛋白从“类别判定”向“程度刻画”的转变提供了可操作路径。

总体而言，PIC 的研究现状表明：在无需显式结构或网络特征的前提下，基于大
规模预训练序列模型的微调策略能够实现高精度的蛋白必需性识别，并具备跨物
种与跨细胞层级扩展能力。这确立了 PLM 驱动范式在必需蛋白预测领域中的可行
性与稳定性。

然而，PIC 驱动模型所定义的“必需性”主要基于细胞存活或增殖能力的依赖
性评估，其核心关注点在于维持基本生命活动所需的蛋白集合。对于高度功能特
化且状态依赖性显著的免疫细胞而言，蛋白的重要性不仅体现在生存层面，更体
现在炎症响应、吞噬活化及代谢重编程等功能维度。因此，如何在保持序列驱动
建模范式一致的前提下，识别特定免疫细胞中的功能核心蛋白，仍然是尚未系统
解决的问题。

基于此，本研究在继承 PLM 微调框架的基础上，将其应用于中性粒细胞特异情
境，构建免疫细胞功能依赖蛋白识别流程，并以模型预测结果为线索提出潜在靶
点蛋白。

\section{问题的提出与研究设计}

\subsection{研究问题与总体思路}

近年来，基于预训练 PLM 的蛋白质必需性预测方法取得显著进展。Protein Importance Calculator（PIC）通过整合人群遗传约束数据与细胞系 CRISPR 筛选数据，构建了多层级蛋白必需性预测框架，并提出 PES 用于量化蛋白在不同生物层面的重要程度\cite{kang2024comprehensive}。该方法显著提升了蛋白必需性预测性能，为从序列层面识别功能关键蛋白提供了可靠工具。

然而，现有研究多集中于泛人类或整体细胞层面的核心必需蛋白识别，其应用目标主要围绕生命维持核心蛋白或肿瘤细胞依赖蛋白的筛选。对于特定免疫细胞群体，尤其是中性粒细胞这一高度动态、强应激的效应细胞，其潜在功能调控靶点尚未通过系统化计算模型进行挖掘。

中性粒细胞在吞噬体酸化、呼吸爆发及强氧化应激环境中执行免疫功能\cite{arocacrevilln2024neutrophils}。其中，急性炎症应答中以氧化爆发增强和颗粒/溶酶体功能负荷升高为特征的活化中性粒细胞状态，即氧化爆发-颗粒/溶酶体高负荷活化状态，尤其依赖维持酸化动力学、膜稳态和蛋白降解效率的关键分子。然而，目前尚缺乏基于深度学习必需性预测模型、围绕这一状态系统筛选中性粒细胞候选靶点蛋白的研究。

因此，本研究的核心目标在于：基于深度学习蛋白必需性预测模型，系统筛选在免疫情境中具有条件性依赖、并与氧化爆发-颗粒/溶酶体高负荷活化状态相关的候选靶点蛋白，并建立从计算预测到机制验证的研究流程。

围绕该目标，本研究形成“模型预测—候选筛选—机制验证”的总体思路。首先，利用预训练 PLM 提取序列表征，构建面向免疫情境的必需性预测模型。本文所称免疫情境，是指免疫效应执行过程中由氧化应激增强、胞内运输与膜重塑活跃以及颗粒/溶酶体功能负荷升高共同构成的高应激功能情境。鉴于本文始终以中性粒细胞为研究对象，该情境在生物学解释和实验验证层面进一步具体化为急性炎症应答中以氧化爆发增强和颗粒/溶酶体功能负荷升高为特征的活化中性粒细胞状态，即氧化爆发-颗粒/溶酶体高负荷活化状态。由此，Human-level 与 Immune-level 的比较对应基础生存依赖与条件性依赖的层级比较，中性粒细胞则是本文进行生物学聚焦的具体研究对象。在实验设计上，HL-60 细胞被选作中性粒细胞的体外替代体系。已有研究表明，PMA 刺激可在 HL-60 体系中诱导 NADPH 氧化酶相关 ROS 生成，并在中性粒细胞样 HL-60 体系中更明确地触发氧化应激应答 \cite{muranaka2005mechanism,kuwabara2015nadph}。因此，本文在 HL-60 替代体系中施加 PMA 刺激，用以模拟氧化爆发-颗粒/溶酶体高负荷活化状态，从而近似再现氧化爆发增强及溶酶体/酸化负荷升高等关键功能特征。其次，以 PES 为核心量化指标筛选高置信候选蛋白，并结合功能富集与结构特征分析评估其生物学合理性；最后，选取代表性候选蛋白开展分子结构模拟与体外实验验证，形成从计算筛选到机制解析的系统验证流程。
\begin{figure}
    \centering
    \includegraphics[width=1\linewidth]{figures/reflow.png}
    \caption[研究流程]{研究流程。该流程概述了本研究从免疫情境蛋白必需性预测模型构建、候选蛋白筛选，到围绕氧化爆发-颗粒/溶酶体高负荷活化状态开展结构机制解析与实验验证的整体研究路径，强调“模型预测—生物学筛选—机制验证”的系统设计。}
    \label{fig:research_workflow}
\end{figure}


\subsection{技术路线与论文结构安排}

在技术实现层面，本研究采用基于 T5 架构的预训练蛋白语言模型 ProtT5 提取序列嵌入特征，并构建基于注意力机制（attention mechanism, Attention）和多层感知机（multilayer perceptron, MLP）的多模块深度学习架构进行蛋白必需性预测。模型整体延续 Embedding--Attention--Prediction 的设计思路，并通过消融实验对 Attention 模块进行优化。实验结果显示，单头 Attention 结构在精准率-召回率曲线下面积（area under the precision-recall curve, AUPRC）指标上优于多头结构，提示在必需性预测任务中，集中式注意力机制更有利于关键序列区域的识别与表达。

模型性能通过 AUROC 与 AUPRC 等指标进行系统评估，并采用交叉验证保证模型稳健性。在模型预测结果基础上，本研究依据 PES 评分对蛋白进行排序筛选，构建高置信候选集合。通过比较不同模型版本中评分稳定的蛋白，识别在免疫情境下持续高评分的核心候选蛋白。

随后，对候选蛋白开展理化与结构特征分析，包括等电点分布、表面电势特征、疏水性指数等指标，以评估其在免疫高应激环境中的潜在适应性。
在此基础上，选取代表性候选蛋白，开展结构机制解析与验证。

全文结构安排如下：第一章为引言，系统阐述蛋白必需性研究背景与方法学发展现状；第二章介绍免疫情境蛋白必需性预测模型的构建与优化；第三章基于模型结果筛选中性粒细胞相关候选蛋白并开展结构与功能分析；第四章围绕氧化爆发-颗粒/溶酶体高负荷活化状态下的核心候选蛋白开展分子机制解析与实验验证；第五章总结研究成果并展望未来应用方向。

本章系统梳理了蛋白质必需性鉴定方法的技术演进路径，并评述了 PLM 在序列语义建模中的方法学突破。通过对中性粒细胞生物学特性及现有筛选技术局限的分析，可以明确当前研究在细胞类型特异性必需蛋白识别方面仍存在明显空白。

基于上述问题界定，后续章节将围绕免疫情境中的蛋白必需性预测展开模型构建与验证，并在生物学解释与实验层面聚焦氧化爆发-颗粒/溶酶体高负荷活化状态，利用深度学习方法从海量序列信息中挖掘维持中性粒细胞功能稳态的关键分子，预期能够发现新的免疫调控靶点蛋白。
